\documentclass[10pt,a4paper]{amsart}
\usepackage[margin=19mm]{geometry}
\usepackage{amsmath,amssymb,mathtools}
\usepackage[scr=rsfs,cal=euler]{mathalfa}
\usepackage{xfrac}
\usepackage[unicode,hidelinks,bookmarks=false]{hyperref}
\newcommand{\ie}{i.e.\ }
\newcommand{\cf}{cf.\ }
\newcommand{\resp}{respectively\ }
\newcommand{\Subsection}{\subsection}
\providecommand{\nobreakdash}{\nobreak\hbox{-}}
\setlength{\emergencystretch}{2em}
\multlinegap0pt
\usepackage{dsfont}
\usepackage{bm}
\usepackage{amssymb}
\usepackage{mathtools}
\usepackage{tikz}
\usepackage{algorithm}\usepackage{algpseudocode}
\usepackage{csquotes}
\newtheorem{lemma}{Lemma}
\newtheorem{theorem}{Theorem}
\renewcommand{\Im}[1]{\mathrm{Im}(#1)}
\renewcommand{\Re}[1]{\mathrm{Re}(#1)}
\renewcommand{\d}{\,\mathrm{d}}
\newcommand{\new}[1]{#1} 
\newcommand{\zeps}[3]{\new{Z_{#1} (#2, #3)}}
\title{Computation and Properties of the Epstein Zeta Function with Applications to Quantum Systems}
\author{Andreas A. Buchheit, Jonathan Busse, Ruben Gutendorf}
\date{Source: arXiv:2412.16317v2 (CC BY 4.0)}
\begin{document}
\maketitle

\noindent\emph{Source and licence.} Computation and Properties of the Epstein Zeta Function with Applications to Quantum Systems. Version arXiv:2412.16317v2 (CC BY 4.0), distributed by arXiv under the Creative Commons Attribution 4.0 International licence, http://creativecommons.org/licenses/by/4.0/, as recorded in the arXiv licence metadata for this exact version and shown on the abstract page https://arxiv.org/abs/2412.16317v2. Full licence notice: \texttt{licenses/arxiv-2412.16317-CC-BY-4.0.txt}.\par\smallskip

\noindent\textit{Editorial scope.} Counted unit: the article's theorem theorem:trunc (source0.tex lines 1183--1226). The article's complete original proof of it is delivered with it (source0.tex lines 1964--2044, one \texttt{proof} environment of 81 lines, which the article places away from the statement). No mathematics is written, repaired or re-derived: the statement and the proof are the article's own lines, in the article's own notation, and no retained line is substituted or re-flowed.  Retained uncounted as the source-local context the proof needs: lines 428--460; lines 1866--1873; lines 1886--1900; lines 1942--1946.  Nothing was omitted from any retained span, so the package carries no omission record.  No retained line contains a citation, so the wrapper carries no bibliography and no bibliography entry.  No retained line contains a figure environment, an image-inclusion command or drawing code, so no image is delivered and the package declares no asset dependency.  Editorial formatting only, in addition: an AMS \texttt{amsart} preamble that restates the article's own declarations and macros used by the retained lines, namely the environments \texttt{lemma}, \texttt{theorem} on separate counters, together with the standard packages those lines need.  The retained environments print in this document as Lemma 1, Theorem 1 in order of appearance, whereas the article prints the same items with the numbering of its own full document.  The complete native source of the article as distributed in the arXiv e-print for this exact version is bundled unchanged as \texttt{sources/170506/source0.tex}.
\begin{lemma}[Properties of Crandall functions]
\label{lem:propCrandall}
Let $\nu \in\mathds C$ and $\bm z\in D$ with the complex sector
$$
D=\{\bm u\in\mathds C^d:|\Re{\bm u}|>|\Im{\bm u}|\},
$$
where real and imaginary parts are applied component-wise and where $\vert \bm \cdot \vert$ denotes the Euclidean norm.
\begin{enumerate}
\item (Holomorphy) The lower Crandall function $g_{\nu}(\bm z)$ 
\new{can be jointly holomorphically extended to} 
$(\nu,\bm z)\in\mathds C\setminus(-2\mathds N_0)\times \mathds C^d$. The upper Crandall function $G_{\nu}(\bm z)$ can be extended to a jointly holomorphic function in $(\nu,\bm z)\in \mathds C \times D$. Finally,
$G_{\nu}(\bm 0)=-2/\nu$ is holomorphic in $\nu\in\mathds C\setminus\{0\}$
with a simple pole at $\nu=0$.
\item (Fundamental relation)
For any $\lambda>0$, it holds
\[
( \bm z^2 )^{-\nu/2} = \frac{(\pi/\lambda^2 )^{\nu/2}}{\Gamma(\nu/2)} \Big(G_\nu(\bm z/\lambda)+g_\nu(\bm z/\lambda)\Big).
\]
Further, for $\mathrm{Re}(\nu)<0$, we have $\lim_{\alpha \to 0_+} G_\nu(\alpha\bm z)=G_\nu(\bm 0)=-2/\nu$. For all $\nu \in \mathds C$, we have
\[
\frac{1}{\Gamma(\nu/2)}(G_\nu(\bm 0)+g_\nu(\bm 0)) = 0,
\]
\new{where division by the Gamma function removes the poles at $\nu\in -2 \mathds N_0$.}
\item (Uniform bound)
Define
\[
u=\pi \mathrm{Re}(\bm z^2),\quad v=\max\{0,\mathrm{Re}(\nu/2)-1\}.\] 
We then have
$$
|G_\nu(\bm z)| \leq \frac{e^{-u}}{u-v}, \quad u>v.
$$
\end{enumerate}
\end{lemma}

\begin{lemma}
\label{lem:z-is-enough-revisited}
Let $\Lambda=A\mathds Z^d$ for $A\in\mathds R^{d\times d}$ regular.
Then, for any set $M\subset\mathds R^d\setminus\{\bm 0\}$,
and $\bm v\in\mathds R^d$, we have
\[\sum_{\bm z \in \Lambda \cap M} G_\nu(c(\bm z+\bm v)) \leq \sum_{\bm z \in \mathds{Z}^d \cap A^{-1}M} G_\nu\left(\frac{c(\bm z +A^{-1}\bm v)}{\Vert A^{-1}\Vert} \right),\]
with $\|\cdot\|$ the spectral norm.
\end{lemma}

\begin{lemma}
\label{lem:shells}
Let $d\in\mathds N_+$, let $r>0$, and let
$$
N_{d}(r)=\sup_{\bm v\in \mathds R^d}\big(\# 
\{\bm z\in\mathds Z^d: |\bm z+\bm v|\le r\}\big),
$$
where the cardinality $\#$ denotes the number of elements in a set. 
Then, for any $r>\sqrt{d}$ and $0<\varepsilon\le (r-\sqrt{d})/2$, it holds that
$$
N_d(r+\varepsilon) 
\le 
\frac{(3/2)^d\pi^{d/2}}{\Gamma(d/2+1)}r^d.
$$
\end{lemma}

\begin{lemma}
\label{lem:g-int-revisited}
Let $t \in \mathds{R}$ and let $s \in \mathds{R} \setminus \{t\}$. Then, for any $r > 0$, we have
\[\int_r^\infty u^{t}G_s(u)\frac{\d u}{u} = r^{t} \frac{G_{t}(r) - G_s(r)}{t-s}.\]
\end{lemma}

\begin{theorem}[Remainder estimate]
  \label{theorem:trunc}
Let $\Lambda=A\mathds Z^{d\times d}$, for $A\in\mathds R^{d\times d}$ regular and $\det(A)=1$. Let $\bm x, \bm y \in \mathds{R}^d$ and let
$\nu\in I$ with $I\subset\mathds R$ compact so that $\nu\neq d$ if $\bm y\in\Lambda^*$.
Let $\mathcal{R}_{\Lambda}(r)$ denote the remainder obtained by truncation at $r>0$ in Crandall's representation
\begin{align*}
\zeps{\Lambda,\nu}{\bm{x}}{\bm{y}}
=&\frac{\pi ^{\nu/2}}{\Gamma(\nu/2)}\Bigg[\sum_{\substack{\bm z\in\Lambda\\ |\bm z - \bm x|\le r}}
G_{\nu}(\bm{z} -\bm x)e^{-2\pi i\bm{y}\cdot\bm{z}}
\\
&
\quad
+\frac{1}{V_{\Lambda}}
\sum_{\substack{\bm k\in\Lambda^{\ast}\\ |\bm k + \bm y|\le r}}G_{d-\nu}
(\bm{k} + \bm y) e^{2\pi i\bm{x}\cdot (\bm{k} + \bm y)}
\Bigg]+\mathcal R_\Lambda(r).
\end{align*}
Then, for any $\varepsilon>0$ and $r>(1+2\varepsilon)\sqrt{d}\,\kappa(A)$ with the condition number $\kappa(A)=\|A\|\|A^{-1}\|$, where $\|\cdot\|$ denotes the spectral norm,
it holds that
$$
|\mathcal R_\Lambda(r)|<\kappa(A)^{d+1} \sup_{\nu \in I}
c_\nu
\Big(
R_{\nu}\big(r/\kappa(A)\big)
+
R_{d-\nu}\big(r/\kappa(A)\big)
\bigg)
$$
with the prefactor \[c_\nu = \frac{(3/2)^d\pi^{(\nu+d)/2}}{\Gamma(d/2+1)|\Gamma(\nu/2)|},\]
and with the function
$$
R_{\nu}(r)
=\frac {r^{d+1}}\varepsilon
\Big(
\frac{G_{d+1}(r-\varepsilon) - 
G_{\nu}(r-\varepsilon)}{d+1-\nu}\Big),
$$
where the limit $\nu\to d+1$ is well-defined.
In particular, 
\[
\mathcal R_\Lambda(\kappa(A)r)
\le \kappa(A)^{d+1}
\mathcal R_{\mathds Z^d}(r).\]
\end{theorem}

\begin{proof}[Proof of Theorem \ref{theorem:trunc}]
By 
Lemma \ref{lem:z-is-enough-revisited} the absolute value of the remainder is bounded by
$$
    \frac{\pi^{\nu/2}}{|\Gamma(\nu/2)|}\Bigg[\sum_{\substack{\bm z\in\mathds Z^d\\ |\bm z - A^{-1}\bm x|> r/\|A\|}}\!\!
    G_{\nu}\Big(\frac{\bm{z} -A^{-1}\bm x}{\|A^{-1}\|}\Big)
+\frac{1}{V_{\Lambda}}
\!\!\!
\sum_{\substack{\bm k\in\mathds Z^d\\ |\bm k + A^T\bm y|> r/\|A^{-1}\|}}
\!\!\!\!\!
G_{d-\nu}
\Big(\frac{\bm{k} + A^T\bm y}{\|A\|}\Big) 
\Bigg]
$$
where we enlarged the summation range using $A^{-1}(\mathds R^d\setminus B_r)\subseteq \mathds R^d\setminus B_{r/\Vert A\Vert}$ and
$A^{T}(\mathds R^d\setminus B_r)\subseteq \mathds R^d\setminus B_{r/\Vert A^{-T}\Vert}$ with $\Vert A^{-T}\Vert=\Vert A^{-1}\Vert$.


We now subdivide the sums into sums over spherical shells of width 
$\varepsilon'>0$ as follows
$$\sum_{\substack{\bm k\in \mathds Z^d\\ |\bm k + \bm v|> r'}}G_{s}
(c(\bm{k} + \bm v)) = \sum_{n = 0}^\infty \sum_{\substack{\bm k\in \mathds Z^d\\ |\bm k + \bm v|> r'+n\varepsilon'\\|\bm k + \bm v|\le  r'+(n+1)\varepsilon' }} G_{s}
(c(\bm{k} + \bm v)) 
$$
for
$r',c>0$, $\bm v\in\mathds R^d$, and $s\in \mathds C$.
We then use the strict monotonic decrease and rotational symmetry of $G_s(\cdot)$ to bound the value of $G_s$ by its maximum within the shell. We further bound the number of lattice points within the shell by the number of points within the closed ball of radius $r'+(n+1)\varepsilon'$, with the bound provided in 
Lemma \ref{lem:shells}. We obtain for 
$0<\varepsilon'\le (r'-\sqrt{d})/2$,
$$
\sum_{\substack{\bm k\in \mathds Z^d\\ |\bm k + \bm v|> r'}}G_{s}
(c(\bm{k} + \bm v))
\le
C\sum_{n=0}^{\infty}
(r'+n\varepsilon')^d
G_{s}
(c(r'+n\varepsilon')),
$$
with $C=(3/2)^d\pi^{d/2}/\Gamma(d/2+1)$.
By monotonicity,
this expression is bounded by the associated integral
$$
C\int_{-1}^{\infty}
(r'+n\varepsilon')^d
G_{s}
(c(r'+n\varepsilon'))\,\mathrm d n.
$$
Applying 
Lemma \ref{lem:g-int-revisited}, we find
$$
C
\frac 1{\varepsilon'}\Big[(r'-\varepsilon')^{d+1}
\frac{G_{d+1}(c(r'-\varepsilon')) - G_s(c(r'-\varepsilon'))}{d+1-s}
\Big]
$$
where the limit $s\to d+1$  is well-defined, since $G_{s}(u)$, $u>0$,
is holomorphic in $s\in\mathds C$ by 
Lemma \ref{lem:propCrandall}.
Letting
 $r'=r/\|A\|$, $c=1/\|A^{-1}\|$
 for the first sum, and $r'=r/\|A^{-1}\|$, $c=1/\|A\|$
 as well as $\tilde\varepsilon=c \,\varepsilon'$
 leads to the bound
$$
\kappa(A)^{d+1}C
\frac 1{\tilde\varepsilon}\Big[
\Big(\frac r{\kappa(A)}\Big)^{d+1}
\frac{G_{d+1}(r/\kappa(A)-\tilde\varepsilon) - G_s(r/\kappa(A)-\tilde\varepsilon)}{d+1-s}
\Big]
$$
where we used that $c\le 1$ since both $\|A\|\ge 1$ and $\|A^{-1}\|\ge 1$ due to $\det(A)=1$.
The bound above holds
for every
$$
0<\tilde\varepsilon\le 
(r/\kappa(A)-\sqrt{d}/\max\{\|A\|,\|A^{-1}\|\})/2.
$$
The right-hand side of the expression above is, by the constraint on $r$, larger than $\varepsilon$.
Thus, we may choose $\tilde{\varepsilon}=\varepsilon$.
Including the $\nu$-dependent prefactor yields the statement.
\end{proof}

\end{document}
